\documentclass[11pt,a4paper]{article}
\usepackage[T1]{fontenc}
\usepackage[utf8]{inputenc}
\usepackage{lmodern,microtype}
\usepackage[margin=25mm,headheight=15pt]{geometry}
\usepackage{amsmath,amssymb,amsthm,mathtools,mathrsfs}
\usepackage{booktabs,longtable,array,enumitem}
\usepackage[dvipsnames]{xcolor}
\usepackage{fancyhdr,titlesec,xurl}
\usepackage[colorlinks=true,linkcolor=MidnightBlue,citecolor=MidnightBlue,urlcolor=MidnightBlue]{hyperref}
\usepackage{bookmark}
\makeatletter
\renewcommand\l@subsection{\@dottedtocline{2}{1.5em}{3.1em}}
\makeatother
\definecolor{ink}{RGB}{24,49,69}
\titleformat{\section}{\large\bfseries\color{ink}}{\thesection}{0.65em}{}
\titleformat{\subsection}{\normalsize\bfseries\color{ink}}{\thesubsection}{0.65em}{}
\pagestyle{fancy}\fancyhf{}
\fancyhead[L]{\small Amplitude--slope compensation in feedback transseries}
\fancyhead[R]{\small Research article}
\fancyfoot[C]{\thepage}
\renewcommand{\headrulewidth}{0.3pt}
\setlength{\parskip}{3pt}
\setlength{\emergencystretch}{3em}
\setlist{itemsep=4pt,topsep=5pt,leftmargin=1.8em}
\numberwithin{equation}{section}
\newtheorem{theorem}{Theorem}[section]
\newtheorem{proposition}[theorem]{Proposition}
\newtheorem{lemma}[theorem]{Lemma}
\newtheorem{corollary}[theorem]{Corollary}
\theoremstyle{definition}
\newtheorem{definition}[theorem]{Definition}
\newtheorem{example}[theorem]{Example}
\theoremstyle{remark}
\newtheorem{remark}[theorem]{Remark}
\newcommand{\R}{\mathbb R}
\newcommand{\C}{\mathbb C}
\newcommand{\N}{\mathbb N}
\newcommand{\U}{\widehat U}
\newcommand{\Q}{\widehat Q}
\newcommand{\Gev}{\mathcal G}
\newcommand{\B}{\mathcal B_0}
\newcommand{\e}{\mathrm e}
\newcommand{\dd}{\mathrm d}
\newcommand{\Ree}{\operatorname{Re}}
\newcommand{\Imm}{\operatorname{Im}}
\newcommand{\dist}{\operatorname{dist}}
\newcommand{\inv}{\langle-1\rangle}
\newcommand{\coeff}[2]{[#1^{#2}]}
\newcommand{\calA}{\mathcal A}
\newcommand{\calD}{\mathcal D}
\hypersetup{pdftitle={Amplitude-Slope Compensation in Feedback Transseries},pdfauthor={Research draft prepared with ChatGPT for Vladimir Reshetnikov},pdfsubject={Weighted exponential feedback: convergence, Gevrey classes, fine summability, exact type, and condensation}}
\begin{document}
\begin{center}
{\LARGE\bfseries\color{ink} Amplitude--Slope Compensation\par in Feedback Transseries\par}
\vspace{0.8em}
{\large Sharp regularity, negative-ray summation,\par and damping-induced condensation\par}
\vspace{1em}
Research article prepared for Vladimir Reshetnikov\\
29 September 2026
\end{center}

\begin{abstract}
Consider the formal equation
\[
 \U(q)=\sum_{j\ge1}c_jq^j\exp(\lambda_j\U(q)),\qquad
 c_1=1,\quad 0<c_j\le1,\quad\lambda_j\ge0,
\]
and let $d_j=\log(1/c_j)$. We prove that its solution converges precisely
when $\lambda_j=O(j+d_j)$. For every $s>0$, the solution and its
compositional inverse are Gevrey--$s$ precisely when
\[
 \lambda_j^{1/(s+1)}=O\bigl(j\log(j+1)+d_j\bigr).
\]
The same condition characterizes a uniform Gevrey--$s$ expansion of the
literal inverse on a closed left half-disc. At $s=1$ it is also equivalent
to fine Borel--Laplace summability of both series in the negative direction;
the sums are inverse functions and satisfy the original convergent kernel.
Under strong damping, $d_j/(j\log(j+1))\to\infty$, the exact
factorial-normalized type is
\[
 T_s(\U)=(s+1)^{s+1}\limsup_{j\to\infty}
                          \frac{\lambda_j}{d_j^{s+1}}.
\]
For $d_j\sim b j^\beta$, $\lambda_j\sim a j^p$, $p>\beta>1$, the type
is an actual coefficient-root limit, and the dominant action has size
$((p/\beta)n/b)^{1/\beta}$ in probability. These results answer the
repository's weighted-amplitude question for positive, exponentially
bounded amplitudes. The package includes exact finite verification,
a rational inverse certificate, numerical diagnostics, and further
research questions.
\end{abstract}

\noindent\textbf{Status and scope.} These are conventional proofs, not
Lean-verified statements. Global publication priority has not been established.
The contribution claimed is a precise extension of the inspected ProveIt
research, not a solution of an unrelated named conjecture. Fine summability
means a fixed-width neighborhood of one Borel ray; angular summability and
resurgence are not asserted.

\clearpage
\setcounter{tocdepth}{2}
\tableofcontents
\clearpage

\section{The question addressed and the contribution}
\label{sec:context}

\subsection{Repository provenance}
The source snapshot inspected for this article is
\begin{center}\small\ttfamily
ebd8344bca77d8352cd745b7df374618a290a029.
\end{center}
The relevant project is \texttt{Analysis/Transseries/}. Its directory
README describes the recent research packages as unmerged and not
formalized \cite{project}. We use their mathematical questions and
methods as starting points, not their presence in a repository as a
certificate of correctness.

Two specific source passages motivate the work. The regularity article
asks for an amplitude-dependent analogue of its unit-amplitude results
\cite[lines 1470--1483]{regularity}. The negative-ray article asks for a
sharp joint condition on amplitudes and slopes, including very rapidly
decaying amplitudes \cite[section ``Weighted actions and joint
amplitude--slope thresholds'']{negative}. The latter source proves a
unit-amplitude fine-summation theorem. The present paper treats an
arbitrary positive amplitude sequence with an exponential upper bound.
It does not change the action exponents from the integers $j$.

The focused review included the project and document-directory READMEs,
the relevant regularity questions, and the negative-ray article's
statements, inverse representation, continuation argument, and open
questions. It was not an exhaustive audit of the canonical volume or of
every concurrent research arrival. Reproducible source pointers are
recorded in \texttt{SOURCE\_NOTES.md}.

\subsection{Why amplitudes change the problem}
For unit amplitudes, a large slope $\lambda_j$ can be used repeatedly
through the exponential expansion without paying an initial suppression
cost. A small amplitude $c_j$ changes the balance. Entering action $j$
now costs $d_j=-\log c_j$ in the logarithm of a positive coefficient.
The cost can be larger than the combinatorial penalty associated with
its action index. It can even compensate for slopes growing faster than
any polynomial.

Three separate questions must therefore be distinguished. First, does the
formal solution have a nonzero ordinary convergence radius? Second, which
factorial growth bounds do its coefficients satisfy? Third, does an
analytic realization recover those coefficients with uniform remainders
and a canonical directional summation rule? Our results address all
three, but the analytic remainder statement for arbitrary $s$ is made
for the inverse. Fine summability of the forward series follows at
$s=1$ from a separate nonlinear argument.

\subsection{What is new relative to the inspected sources}
The principal additions are the joint necessary-and-sufficient criteria,
the weighted uniform remainder theorem, and the exact type identity in
the strong-damping regime. The stretched-exponential specialization
exhibits a dominant-action scale $n^{1/\beta}$ and a nonzero critical
factorial type that differ from the unit-amplitude behavior. No
monotonicity or regular variation is needed for the joint criteria or
for the strong-damping \emph{limsup} identity.

The algebraic tools are classical. Lagrange inversion is used in its
standard coefficient form \cite{gessel}. The Borel kernels are modified
Bessel functions, with an elementary integral representation
\cite{dlmf}. The distinction between fine and angular summability, and
the convolution approach to nonlinear operations, are standard
\cite{sauzin}. The new input is the quantitative use of the amplitude
cost throughout those constructions. We prove the specialized estimates
needed here rather than appealing to a generic analytic implicit-function
theorem for a kernel that may not be analytic at the origin.

\section{Definitions and the main results}
\label{sec:main}

\subsection{Normalization and formal series}
Throughout the main proofs assume
\begin{equation}
 c_1=1,\qquad 0<c_j\le1,\qquad \lambda_j\ge0,\qquad
 d_j=-\log c_j\ge0.
 \label{eq:normalization}
\end{equation}
All individual slopes are finite. Define
\begin{equation}
 \Phi(q,u)=\sum_{j\ge1}c_jq^j\e^{\lambda_j u},\qquad
 \U=\Phi(q,\U),\qquad
 \U(q)=\sum_{n\ge1}u_nq^n,
 \quad \Q=\U^{\inv}=\sum_{n\ge1}v_nu^n.
 \label{eq:model}
\end{equation}
Here $u$ is the argument of $\Q$, not a coefficient. Since $u_1=1$, the
formal compositional inverse exists. For any fixed degree, every sum
used in the formal equation is finite.

\begin{definition}[Gevrey order and type]
For $s>0$, a formal series $f(z)=\sum f_nz^n$ belongs to $\Gev_s$ if
$|f_n|\le CH^n(n!)^s$ for some $C,H>0$. Its exact
factorial-normalized type is
\begin{equation}
 T_s(f)=\limsup_{n\to\infty}
                  \left(\frac{|f_n|}{(n!)^s}\right)^{1/n}
                  \in[0,\infty].
 \label{eq:type}
\end{equation}
Thus $f\in\Gev_s$ if and only if $T_s(f)<\infty$. The convention uses
$(n!)^s$, not $\Gamma(1+sn)$; the normalization matters for exact constants.
\end{definition}

Set
\begin{equation}
 L_j=j\log(j+1),\qquad H_j=L_j+d_j.
 \label{eq:Hj}
\end{equation}
A uniform Gevrey expansion of $Q$ on a closed left half-disc means that,
for some $r_0,C,A>0$ and every $N\ge1$,
\begin{equation}
 \left|Q(u)-\sum_{n=1}^{N-1}v_nu^n\right|
 \le CA^N(N!)^s|u|^N,
 \qquad \Ree u\le0,\quad |u|\le r_0.
 \label{eq:uniform-G}
\end{equation}
Boundary values on the imaginary axis are furnished by the convergent
block construction, not by analytic continuation across that axis.

\begin{theorem}[Joint convergence and regularity classification]
\label{thm:classification}
Under \eqref{eq:normalization}, the formal solution converges in a
neighborhood of zero if and only if
\begin{equation}
 \sup_{j\ge1}\frac{\lambda_j}{j+d_j}<\infty.
 \label{eq:analytic-criterion}
\end{equation}
For every fixed $s>0$, the following assertions are equivalent:
\begin{enumerate}[label=(\roman*)]
 \item $\displaystyle\sup_j\lambda_j^{1/(s+1)}/H_j<\infty$;
 \item $\U\in\Gev_s$;
 \item $\Q\in\Gev_s$;
 \item the literal inverse branch $u=\Phi(Q(u),u)$ satisfies
       \eqref{eq:uniform-G} on a sufficiently small closed left half-disc.
\end{enumerate}
The literal inverse branch exists on the left independently of these
growth conditions. No monotonicity hypothesis is required.
\end{theorem}

\begin{definition}[Fine Borel summability]
Let
\[
 \B f(\zeta)=\sum_{n\ge0}\frac{f_n}{n!}\zeta^n,
 \qquad
 \Omega^-_\delta=\{\zeta:\dist(\zeta,(-\infty,0])<\delta\}.
\]
A formal series is finely Borel summable in direction $\pi$ when its
Borel germ extends holomorphically to some $\Omega^-_\delta$ and is bounded
there by $C\e^{A|\zeta|}$. Its sum at $u=-r$ is
\begin{equation}
 \mathcal S_\pi f(-r)=\frac1r\int_0^\infty
                 \e^{-t/r}\B f(-t)\,\dd t.
 \label{eq:laplace}
\end{equation}
\end{definition}

\begin{theorem}[Sharp weighted negative-ray summability]
\label{thm:fine}
Under \eqref{eq:normalization}, each of the four conditions of
Theorem~\ref{thm:classification} at $s=1$ is also equivalent to each of
\[
 \U\text{ is finely summable in direction }\pi,
 \qquad
 \Q\text{ is finely summable in direction }\pi.
\]
The sums are compositional inverses near the negative real axis. They
satisfy the literal kernel equation on their sufficiently small common
left-side domain, with an absolutely convergent sum over $j$.
\end{theorem}

\begin{theorem}[Exact type under strong amplitude damping]
\label{thm:type}
Suppose additionally that
\begin{equation}
 \frac{d_j}{j\log(j+1)}\longrightarrow\infty.
 \label{eq:strong}
\end{equation}
Then for every $s>0$, with values in $[0,\infty]$,
\begin{equation}
 \boxed{\quad
 T_s(\U)=(s+1)^{s+1}
       \limsup_{j\to\infty}\frac{\lambda_j}{d_j^{s+1}}.
 \quad}
 \label{eq:exact-type}
\end{equation}
This is a forward-series type theorem. It does not assert equality of
forward and inverse exact types.
\end{theorem}

\begin{theorem}[Dense power costs: root limit and dominant action]
\label{thm:condensation}
Assume
\begin{equation}
 d_j\sim b j^\beta,\qquad \lambda_j\sim a j^p,
 \qquad a,b>0,\quad p>\beta>1.
 \label{eq:powers}
\end{equation}
Put $r=p/\beta>1$ and $s=r-1$. Then
\begin{equation}
 \lim_{n\to\infty}\left(\frac{u_n}{(n!)^s}\right)^{1/n}
     =a\left(\frac r b\right)^r.
 \label{eq:root-limit}
\end{equation}
Give each ordered composition $(j_1,\ldots,j_k)$ of $n$ probability
proportional to its nonnegative Lagrange weight in
\eqref{eq:ordered-Lagrange}. If $K=k$, $D=\sum_i d_{j_i}$, and
$J_{\max}=\max_i j_i$, then
\begin{equation}
 \frac K n\to1,\qquad \frac D n\to r,\qquad
 \frac{\max_i d_{j_i}}{D}\to1,\qquad
 \frac{J_{\max}}{n^{1/\beta}}\to\left(\frac r b\right)^{1/\beta}
 \label{eq:prob-limits}
\end{equation}
in probability. The ratio involving $D$ is defined arbitrarily when
$D=0$, an event whose probability tends to zero.
\end{theorem}

\subsection{Positive amplitudes with an exponential upper bound}
The normalization $c_j\le1$ does not exclude the usual exponentially
bounded case. Suppose $c_1=1$, $c_j>0$, and
$\sup_{j\ge2}c_j^{1/(j-1)}<\infty$. Choose $R\ge1$ such that
$c_jR^{1-j}\le1$ and set
\begin{equation}
 V(z)=R\U(z/R),\qquad
 \widetilde c_j=c_jR^{1-j},\qquad
 \widetilde\lambda_j=\lambda_j/R.
 \label{eq:rescaling}
\end{equation}
Then $V$ satisfies the normalized model. With
$d_j^+=\max(0,-\log c_j)$, one has
\[
 d_j^+\le\widetilde d_j\le d_j^++(j-1)\log R.
\]
Consequently the analytic and Gevrey criteria hold for the original model
with $d_j$ replaced by $d_j^+$. Fine summability and uniform remainder
properties survive this positive real rescaling. The strong-damping
type formula transforms back to exactly \eqref{eq:exact-type}.

\section{Finite algebra and the necessary estimates}
\label{sec:formal}

\subsection{Existence, positivity, and Lagrange coefficients}
If $V,W\in q\R[[q]]$ agree through degree $n-1$, then
$\Phi(q,V)$ and $\Phi(q,W)$ agree through degree $n$. Iterating from zero
therefore constructs the unique formal solution. All coefficients are
nonnegative, and the coefficient of $q$ is $c_1=1$.

Apply Lagrange inversion to $V=t\Phi(q,V)$ in the auxiliary variable
$t$. Substitution $t=1$ is legitimate coefficientwise because $\Phi$ has
positive $q$-order. The coefficient formula is
\begin{equation}
 u_n=\sum_{k=1}^{n}\frac1{k!}
     \sum_{\substack{j_1+\cdots+j_k=n\\j_i\ge1}}
        \left(\prod_{i=1}^k c_{j_i}\right)
        \left(\sum_{i=1}^k\lambda_{j_i}\right)^{k-1}.
 \label{eq:ordered-Lagrange}
\end{equation}
For $k=1$ the last factor is interpreted as one, even if its base is zero.
Grouping equal parts gives the equivalent finite expression
\begin{equation}
 u_n=\sum_{\sum j m_j=n}
      \frac{\prod_j c_j^{m_j}}{\prod_j m_j!}
            \left(\sum_jm_j\lambda_j\right)^{\sum m_j-1}.
 \label{eq:unordered-Lagrange}
\end{equation}
These formulas are identities in finitely many variables at each degree;
no analytic convergence is being used.

The term with one part $j\ge2$ and $m$ parts equal to one proves
\begin{equation}
 u_{j+m}\ge c_j\frac{(\lambda_j+m\lambda_1)^m}{m!}
              \ge \e^{-d_j}\frac{\lambda_j^m}{m!}.
 \label{eq:one-large}
\end{equation}
Retaining the amplitude in this inequality is essential.

\subsection{Necessity of the joint Gevrey budget}
Suppose $u_n\le CH^n(n!)^s$, enlarging $C,H$ to be at least one. For
$j\ge2$ set $m=\lceil L_j+d_j\rceil$. From \eqref{eq:one-large},
\begin{equation}
 \lambda_j\le C^{1/m}H^{1+j/m}\e^{d_j/m}
                  (m!)^{1/m}((j+m)!)^{s/m}.
 \label{eq:necessary-estimate}
\end{equation}
The factors $C^{1/m}$, $H^{1+j/m}$, and $\e^{d_j/m}$ are bounded.
Also $j+m\le2m$ and
\[
 ((j+m)!)^{s/m}\le(j+m)^s
                \exp\!\left(s\frac j m\log(j+m)\right).
\]
The last exponential is uniformly bounded. Indeed $m\ge j\log(j+1)$,
and $\log(2m)/m$ is decreasing once $m$ is sufficiently large; its
maximum in that range is attained, up to finitely many cases, at the
lower endpoint. Thus $\lambda_j\le C_s m^{s+1}$, which is exactly the
necessary bound in Theorem~\ref{thm:classification}. The finitely many
small indices are absorbed into the constant.

\subsection{Gevrey membership is invariant under formal reversion}
\begin{lemma}
\label{lem:reversion}
For $s>0$ and $f(z)=z+O(z^2)$, one has
$f\in\Gev_s$ if and only if $f^{\inv}\in\Gev_s$.
\end{lemma}
\begin{proof}
Write $f(z)=z(1+h(z))$. The inverse coefficient is
$n^{-1}[z^{n-1}](1+h(z))^{-n}$. A term with $\ell$ factors is indexed by
positive integers $i_1+\cdots+i_\ell=n-1$. Their factorials satisfy
\[
 \prod_{a=1}^{\ell}(i_a+1)!
 \le2^{\ell-1}(n-\ell+1)!\le2^n n!.
\]
To prove this, move units from smaller parts exceeding one to a largest
part; the product cannot decrease, and eventually all but one part equal
one. Raising the estimate to $s$, bounding the binomial coefficient from
the negative-power expansion by $4^n$, and the number of compositions by
$2^n$, gives $|[z^n]f^{\inv}|\le E^n(n!)^s$. The bounds on the input
coefficients contribute only another exponential factor. Apply the same
argument to the inverse for the reverse implication.
\end{proof}
This deliberately proves only class membership. None of the constants
in this lemma are sharp enough to prove an exact type identity.

\section{The exact convergence criterion}
\label{sec:convergence}

Suppose first that $\U$ converges. Choose $q_0>0$ within its radius.
Since every coefficient is nonnegative, Tonelli's theorem applied to the
formal identity gives an equality of nonnegative sums
\[
 \U(q_0)=\sum_j c_jq_0^j\e^{\lambda_j\U(q_0)}<\infty.
\]
Writing $t=\U(q_0)>0$, each summand is bounded by $t$, so
\[
 t\lambda_j\le d_j+j\log(1/q_0)+\log t.
\]
This yields $\lambda_j=O(j+d_j)$.

Conversely, suppose $\lambda_j\le K(j+d_j)$ with $K>0$. On a small
polydisc with $|u|<1/(2K)$,
\begin{equation}
 c_j|q|^j\e^{\lambda_j|u|}
 \le(|q|\e^{1/2})^j\e^{-d_j/2}.
 \label{eq:analytic-majorant}
\end{equation}
The right side is summable when $|q|<\e^{-1/2}$. Thus $\Phi$ is
holomorphic in a neighborhood of $(0,0)$. The analytic implicit-function
theorem applies to $u-\Phi(q,u)$ because its $u$-derivative at the origin
is one. Its Taylor series must be the unique formal solution.

This establishes \eqref{eq:analytic-criterion}. It also shows why the
$j\log j$ term must not be inserted into the convergence criterion.
For example, $c_j=1$ and $\lambda_j=j\log(j+1)$ give a divergent solution
that belongs to every positive Gevrey class. Ordinary convergence is a
strictly different endpoint problem.

\section{The weighted exponential-block inverse}
\label{sec:blocks}

\subsection{Explicit blocks}
For fixed $u$, first invert $q\mapsto\Phi(q,u)$ with an independent
target variable $w$. Lagrange inversion and a negative-binomial expansion
give, after setting $w=u$,
\begin{equation}
 Q(u)=\sum_{k\ge1}u^k
       \sum_{\substack{\mathbf m=(m_1,\ldots,m_\ell)\\
                        m_i\ge1,\ \sum_i m_i=k-1}}
                 \omega_{k,\ell}\e^{-D_{\mathbf m}}
                                    \e^{b_{k,\mathbf m}u},
 \label{eq:blocks}
\end{equation}
where the empty composition is used when $k=1$, and
\begin{align}
 \omega_{k,\ell}&=\frac{(-1)^\ell}{k}\binom{k+\ell-1}{\ell},\label{eq:weight}\\
 D_{\mathbf m}&=\sum_i d_{m_i+1},\label{eq:cost}\\
 b_{k,\mathbf m}&=\sum_i\lambda_{m_i+1}-(k+\ell)\lambda_1.
 \label{eq:frequency}
\end{align}
In particular, the first block is $u\e^{-\lambda_1u}$. Each block starts
at ordinary degree $k$, so \eqref{eq:blocks} also makes sense formally.
It is exactly $\Q$, by uniqueness of formal compositional inversion.

Let $C_k=\sum_{\mathbf m}|\omega_{k,\ell}|$. For $k\ge2$,
\begin{equation}
 C_k=\frac1k\sum_{\ell=1}^{k-1}
           \binom{k+\ell-1}{\ell}\binom{k-2}{\ell-1}
       \le8^k,
 \label{eq:mass}
\end{equation}
and $C_1=1$. The last bound follows directly from
$\binom{k+\ell-1}{\ell}\le4^k$ and the total number $2^{k-2}$ of
compositions of $k-1$. We use this simple bound rather than an optimized
Schr\"oder-number estimate.

\subsection{A literal inverse for every slope sequence}
Every frequency satisfies
\begin{equation}
 b_{k,\mathbf m}\ge-2k\lambda_1.
 \label{eq:negative-frequency}
\end{equation}
Consequently the absolute value of block $k$ on $\Ree u\le0$ is at most
\[
 \bigl(8|u|\e^{2\lambda_1|u|}\bigr)^k.
\]
Choose $r_0>0$ such that
$8r_0\e^{2\lambda_1r_0}<1/3$. The block series converges uniformly on the
closed left half-disc of radius $r_0$, and normally in its interior.
The same calculation with $w$ in place of the factor $u^k$ proves a
uniform convergence radius for the inverse in the independent variable
$w$. Since it agrees near $w=0$ with the analytic inverse of
$q\mapsto\Phi(q,u)$, substitution $w=u$ is valid there. The bound also
keeps $|Q(u)|<1/2$, where the literal kernel converges absolutely.

We have therefore constructed, with no slope-growth assumption,
\begin{equation}
 u=\sum_{j\ge1}c_jQ(u)^j\e^{\lambda_j u},\qquad
 Q(u)=u+O(u^2).
 \label{eq:literal-inverse}
\end{equation}
The latter estimate is uniform on the half-disc. On every smaller closed
sector strictly inside the left half-plane, Cauchy's estimate gives
$Q'(u)=1+O(u)$, so the analytic forward inverse exists near the negative
axis. Uniform Gevrey control is an additional assertion, proved next.

\subsection{A finite block certificate}
Let $Q_{\le K}$ denote the first $K$ blocks. For $u=-r$ with $r\le r_0$,
put $t=8r\e^{2\lambda_1r}<1$. Then
\begin{equation}
 |Q(-r)-Q_{\le K}(-r)|\le\frac{t^{K+1}}{1-t}.
 \label{eq:block-certificate}
\end{equation}
This certificate does not use Gevrey regularity, and it remains valid
for slopes too large for any finite Gevrey class. Only finitely many
amplitudes and slopes are needed to evaluate the first $K$ blocks.

\section{Sufficiency and uniform Gevrey remainders}
\label{sec:sufficiency}

Fix $s>0$ and write $r=s+1>1$. Assume
$\lambda_j^{1/r}\le K(L_j+d_j)$ for all $j$. Constants in this section
may depend on the model and on $s$, but never on a coefficient order or
a particular composition.

\subsection{A weighted block-moment estimate}
In a composition belonging to block $k$, let $D=\sum_i d_{m_i+1}$ and
put $A_k=k\log(k+1)$. Since $\sum_i(m_i+1)\le2k$ and every $m_i+1\le k$,
\[
 \sum_i L_{m_i+1}\le 2k\log(k+1)=2A_k.
\]
For nonnegative numbers and $r\ge1$, $\sum x_i^r\le(\sum x_i)^r$.
Consequently
\[
 \sum_i\lambda_{m_i+1}\le K^r(2A_k+D)^r.
\]
The subtraction in \eqref{eq:frequency} has magnitude at most
$2k\lambda_1$. Increasing a constant also covers block $k=1$, and gives
\begin{equation}
 |b_{k,\mathbf m}|\le C_s(A_k+D)^r.
 \label{eq:weighted-frequency}
\end{equation}
For every integer $m\ge1$,
\begin{align}
 \e^{-D}\frac{|b|^m}{m!}
 &\le \frac{C_s^m}{m!}\e^{-D}(A_k+D)^{rm}\notag\\
 &\le\frac{E_s^m}{m!}
       \left(A_k^{rm}+\left(\frac{rm}{\e}\right)^{rm}\right).
 \label{eq:moment-budget}
\end{align}
Here we used $(x+y)^{rm}\le2^{rm}(x^{rm}+y^{rm})$ and
$\sup_{D\ge0}\e^{-D}D^{rm}=(rm/\e)^{rm}$. The case $m=0$ is bounded by
one. Thus amplitude damping converts a possibly unbounded frequency into
a factorially controlled moment. Bounding frequencies alone would miss
this mechanism.

\begin{lemma}[Factorial budget]
\label{lem:factorial-budget}
For each $s>0$ there is $B_s\ge1$ such that, for $k\ge1$, $m\ge0$, and
$n=k+m$,
\begin{equation}
 \frac{[k\log(k+1)]^{(s+1)m}}{m!}
 +\frac{m^{(s+1)m}}{m!}
 \le B_s^n(n!)^s.
 \label{eq:factorial-budget}
\end{equation}
In the case $m=0$, powers with exponent zero are interpreted as one.
\end{lemma}
\begin{proof}
Write $r=s+1$. The second summand follows at once from
$m!\ge(m/\e)^m$ and $n!\ge(n/\e)^n$, since
$m^{rm}/m!\le\e^m m^{sm}\le\e^{(s+1)n}(n!)^s$.

For the first summand assume $m\ge1$, and use
$\log h!\ge h\log h-h$. Its logarithm after division by $(n!)^s$ is
at most
\[
 rm\log(A_k/n)-sk\log n-m\log(m/n)+m+sn.
\]
The last three terms other than $-sk\log n$ are $O_s(n)$, because
$-x\log x$ is bounded on $[0,1]$. Put
$x=(k/n)\log(n+1)$. Then $A_k/n\le x$, and
$k\log n\ge nx/2$ for $n\ge2$. The remaining expression is at most
\[
 n\left(r\log^+x-\frac{s}{2}x+O_s(1)\right).
\]
The function $r\log^+x-sx/2$ is bounded above on $[0,\infty)$ because
$s>0$. Exponentiation proves the assertion. The finitely many omitted
cases are absorbed into $B_s$.
\end{proof}

\subsection{Coefficient sufficiency}
Extracting the ordinary coefficient of $u^n$ from \eqref{eq:blocks}
uses $k\le n$ and $m=n-k$. Apply \eqref{eq:moment-budget}, then
Lemma~\ref{lem:factorial-budget}, and finally the mass bound $8^k$.
The sum over the $n$ possible blocks is at most $C H^n(n!)^s$.
It follows that $\Q\in\Gev_s$. Lemma~\ref{lem:reversion} gives
$\U\in\Gev_s$. This proves sufficiency for both formal series without
using any asymptotic formula for individual coefficients.

\subsection{The same budget controls an actual analytic remainder}
For every complex $z$ and integer $m\ge1$, the integral form of the
exponential Taylor remainder gives
\begin{equation}
 \left|\e^z-\sum_{h=0}^{m-1}\frac{z^h}{h!}\right|
 \le\frac{|z|^m}{m!}\e^{\max(0,\Ree z)}.
 \label{eq:exp-remainder}
\end{equation}
For $\Ree u\le0$ and a real block frequency $b$, the exponential factor
on the right is at most one if $b\ge0$. If $b<0$, it is at most
$\e^{2k\lambda_1r_0}$. This latter factor is merely exponential in $k$.

In every block $k<N$, truncate $\e^{bu}$ before degree $N-k$.
The resulting finite sum is exactly $\sum_{n<N}v_nu^n$. By
\eqref{eq:exp-remainder} the error in block $k$ is bounded by $|u|^N$
times the same moment in \eqref{eq:moment-budget}, with $m=N-k$,
and an extra fixed exponential in $k$. The factorial budget and mass
bound therefore give a total at most $C H^N(N!)^s|u|^N$.

For the blocks $k\ge N$, choose $r_0$ as in Section~\ref{sec:blocks}.
Their geometric tail is bounded by $C H_0^N|u|^N$. Adding the two
bounds proves \eqref{eq:uniform-G}, including its boundary values.
The case $N=1$ follows directly from the geometric series.

Conversely, if \eqref{eq:uniform-G} holds, its order-$N+1$ version
shows along the negative ray that
\[
 \frac{Q(u)-\sum_{n<N}v_nu^n}{u^N}\longrightarrow v_N.
\]
Its order-$N$ version then gives $|v_N|\le CA^N(N!)^s$.
Thus a uniform expansion implies $\Q\in\Gev_s$, and the necessity
argument of Section~\ref{sec:formal} applies after reversion.
This completes the Gevrey part of Theorem~\ref{thm:classification}.

\begin{remark}[Why the analytic criterion is different]
The proof of Lemma~\ref{lem:factorial-budget} uses $s>0$ in an essential
way. Setting $s=0$ removes the term that offsets $k\log(k+1)$.
Accordingly, the convergence budget is $j+d_j$, not the result of
formally putting $s=0$ in the positive-Gevrey budget. There are divergent
series belonging to every $\Gev_s$, $s>0$.
\end{remark}

\section{Bessel kernels and fine negative-ray summation}
\label{sec:summation}

\subsection{One exact Borel kernel}
For an integer $k\ge1$ and a real number $b$, define the entire function
\begin{equation}
 K_{k,b}(\zeta)=\sum_{m\ge0}
       \frac{b^m\zeta^{k+m}}{m!(k+m)!}
       =\B(u^k\e^{bu})(\zeta).
 \label{eq:bessel-kernel}
\end{equation}
These are Bessel kernels in elementary power-series normalization.
Their relation to the standard $I_k$ and the Poisson integral for $J_k$
can be read from the formulas in \cite{dlmf}. The following direct
version supplies exactly the estimates we need.

Let $\nu_k$ be the probability measure on $[-1,1]$ with density
proportional to $(1-t^2)^{k-1/2}$. Its odd moments vanish and the beta
integral gives
\[
 \int t^{2m}\,\dd\nu_k(t)
 =\frac{(2m)!k!}{4^m m!(k+m)!}.
\]
Expanding the exponential, and integrating on this compact interval,
therefore proves
\begin{equation}
 K_{k,b}(\zeta)=\frac{\zeta^k}{k!}
       \int_{-1}^1\exp(2\sqrt{b\zeta}\,t)\,\dd\nu_k(t).
 \label{eq:bessel-integral}
\end{equation}
The expression is independent of the square-root choice. Put
\[
 h(\zeta)=|\Ree\sqrt\zeta|,
 \qquad \widetilde h(\zeta)=|\Imm\sqrt\zeta|,
 \qquad h(\zeta)^2=\frac{|\zeta|+\Ree\zeta}{2}.
\]
Equation~\eqref{eq:bessel-integral} implies
\begin{equation}
 |K_{k,b}(\zeta)|\le\frac{|\zeta|^k}{k!}
 \begin{cases}
  \exp(2h(\zeta)\sqrt b),&b\ge0,\\
  \exp(2\widetilde h(\zeta)\sqrt{|b|}),&b<0.
 \end{cases}
 \label{eq:bessel-bound}
\end{equation}
In particular, for $t\ge0$ and $b\ge0$, oscillation gives
$|K_{k,b}(-t)|\le t^k/k!$.

For one fixed $b$ and $r>0$, termwise Laplace integration is justified
absolutely and yields
\begin{equation}
 \frac1r\int_0^\infty\e^{-t/r}K_{k,b}(-t)\,\dd t
       =(-r)^k\e^{-br}.
 \label{eq:one-block-laplace}
\end{equation}
With an absolute value inside the integral, its value is at most $r^k$
for $b\ge0$, by the oscillatory bound, and is $r^k\e^{|b|r}$ for $b<0$,
by the positive Taylor series. This stronger distinction will justify an
infinite exchange that termwise absolute Taylor expansion would not.

\subsection{The all-slope ray transform is not automatically a Borel germ}
Write the internal block sums as in \eqref{eq:blocks} and set
\begin{equation}
 \mathfrak B_Q(-t)=\sum_{k\ge1}\sum_{\mathbf m}
       \omega_{k,\ell}\e^{-D_{\mathbf m}}
                         K_{k,b_{k,\mathbf m}}(-t).
 \label{eq:ray-transform}
\end{equation}
The absolute contribution of block $k$ is bounded by
\[
 8^k\frac{t^k}{k!}\exp(2\sqrt{2k\lambda_1t}).
\]
Young's inequality bounds the last exponent by $\eta k+C_\eta t$.
Summing the resulting exponential series proves continuity on the ray
and a bound $C\e^{Ht}$. Moreover,
\eqref{eq:one-block-laplace} bounds the integrated absolute contribution
of block $k$ by $8^kr^k\e^{2k\lambda_1r}$. Hence Fubini's theorem gives
\begin{equation}
 Q(-r)=\frac1r\int_0^\infty
                    \e^{-t/r}\mathfrak B_Q(-t)\,\dd t
 \label{eq:all-slope-recovery}
\end{equation}
for sufficiently small $r$.
This construction works for every slope sequence. Without further growth
control, however, $\mathfrak B_Q$ need not extend to a holomorphic germ
at the origin. Equation~\eqref{eq:all-slope-recovery} alone is not a
Borel-summability theorem.

\subsection{Amplitude damping gives a complex neighborhood}
Now assume the joint criterion with $s=1$.
Equation~\eqref{eq:weighted-frequency} gives, for positive frequencies,
\[
 \sqrt b\le C(A_k+D),\qquad A_k=k\log(k+1).
\]
Combining this with \eqref{eq:bessel-bound} retains the amplitude cost:
\begin{equation}
 \e^{-D}|K_{k,b}(\zeta)|
 \le\frac{|\zeta|^k}{k!}
    \e^{2Ch(\zeta)A_k}\e^{-(1-2Ch(\zeta))D},\qquad b\ge0.
 \label{eq:cost-borel}
\end{equation}
For negative frequencies use $|b|\le2k\lambda_1$ instead. On any compact
set with $2Ch<1$, the absolute mass of block $k$ is consequently at most
\begin{equation}
 \frac{8^k|\zeta|^k}{k!}
 \exp\left(2Ch(\zeta)k\log(k+1)
                 +2\sqrt{2\lambda_1}\widetilde h(\zeta)\sqrt k\right).
 \label{eq:general-borel-majorant}
\end{equation}
Its logarithm is $-(1-\theta)k\log k+O(k)+O(\sqrt k)$ for a fixed
$\theta<1$ on that compact set. The series
\begin{equation}
 B_Q(\zeta)=\sum_{k\ge1}\sum_{\mathbf m}
       \omega_{k,\ell}\e^{-D_{\mathbf m}}
                            K_{k,b_{k,\mathbf m}}(\zeta)
 \label{eq:borel-series}
\end{equation}
therefore converges normally on $\{2Ch<1\}$. Its Taylor coefficient of
order $n$ uses only blocks $k\le n$ and equals $v_n/n!$. Thus it is
indeed the analytic continuation of the Borel germ of $\Q$.

\subsection{Growth on a fixed-width tube}
Choose $\delta>0$ so small that the closure of $\Omega^-_{2\delta}$ is
contained in $\{2Ch<1\}$. For $\zeta=-t+iy$, $t\ge1$, $|y|\le\delta$,
\[
 h(\zeta)\le\frac{\delta}{2\sqrt t},\qquad
 \widetilde h(\zeta)\le\sqrt{t+\delta},\qquad
 |\zeta|\le t+\delta.
\]
The logarithm of \eqref{eq:general-borel-majorant} is bounded by
\begin{equation}
 k\log(8(t+\delta))-\log k!
       +\frac{C_1}{\sqrt t}k\log(k+1)+C_2\sqrt{kt}.
 \label{eq:tube-log}
\end{equation}
If $k\le t^2$, the third term is $o(k)$ as $t\to\infty$ and
$C_2\sqrt{kt}\le k+C_3t$. Summing these terms gives a bound
$C\e^{H_1t}$. If $k>t^2$, then
$\log(t+\delta)\le\tfrac12\log k+o(1)$. For sufficiently large $t$,
$C_1/\sqrt t<1/8$ and $C_2\sqrt{t/k}$ is small. The factorial estimate
$\log k!\ge k\log k-k$ then bounds \eqref{eq:tube-log} by
$-\tfrac14 k\log k$. This tail is uniformly bounded.
The remaining bounded part of the tube is compactly contained in the
normal-convergence domain. It follows that
\begin{equation}
 |B_Q(\zeta)|\le C\e^{H|\zeta|},\qquad\zeta\in\Omega^-_\delta.
 \label{eq:weighted-fine-bound}
\end{equation}
This proves fine summability of $\Q$ and identifies its sum with the
literal $Q$ by \eqref{eq:all-slope-recovery}.

\subsection{Passing from the inverse to the forward series}
Fine summability is closed under compositional inversion for a series
with linear coefficient one. Because fine and angular summability are
different hypotheses, Appendix~\ref{app:inverse} proves the precise
fixed-tube version used here. Its proof is a convolution argument in
the Borel plane, following the classical framework in \cite{sauzin}.
It shows simultaneously that summation commutes with inversion.

Applying that result to $\Q$ proves fine summability of $\U$ and the
identity between the two summed inverse functions. The literal equation
follows from \eqref{eq:literal-inverse}, on the small common domain where
the dependent variable has negative real part. Conversely, fine
summability implies a holomorphic Borel germ and therefore Gevrey--1
coefficients. Theorem~\ref{thm:classification} then gives the joint
budget. This completes Theorem~\ref{thm:fine}.

\section{Exact forward type in the strong-damping regime}
\label{sec:type}

The next argument is independent of the Borel construction. It uses
positivity of the forward Lagrange weights, and it controls their
exponential scale sharply enough to compute an exact type.

\subsection{A universal upper-type lemma}
\begin{lemma}[A power bound on slopes in cost coordinates]
\label{lem:type-upper}
Fix $r>1$, put $s=r-1$, and suppose, for some $a>0$ and $C\ge0$, that
\begin{equation}
 \lambda_j\le a d_j^r+C\qquad(j\ge1).
 \label{eq:cost-power-bound}
\end{equation}
Then $T_s(\U)\le ar^r$. The same upper bound applies to the sum of any
subcollection of the nonnegative Lagrange terms satisfying the analogous
bound $\sum_i\lambda_{j_i}\le aD^r+CK$.
\end{lemma}
\begin{proof}
For a length-$k$ composition of $n$ let $D=\sum_i d_{j_i}$ and
$\Lambda=\sum_i\lambda_{j_i}$. Since
$\sum_i d_{j_i}^r\le D^r$, we have $\Lambda\le aD^r+Ck$.
Increasing $C$ to a positive constant permits the harmless bound
\[
 \frac{\e^{-D}\Lambda^{k-1}}{k!}
 \le\frac1{k!}\e^{-D}(aD^r+Cn)^k.
\]
There are $\binom{n-1}{k-1}$ ordered compositions. Define
\[
 M_{n,k}=\frac1{k!}\sup_{D\ge0}\e^{-D}(aD^r+Cn)^k.
\]
A first, deliberately coarse estimate follows from
$(x+y)^k\le2^k(x^k+y^k)$:
\begin{equation}
 M_{n,k}\le\frac{2^k}{k!}
   \left(a^k(rk)^{rk}\e^{-rk}+(Cn)^k\right).
 \label{eq:coarse-type}
\end{equation}
The logarithm of the first summand is at most
$s k\log k+O(n)$, uniformly for $1\le k\le n$; the second is at most
$\e^{O(n)}$, since $(Cn)^k/k!\le\e^{Cn}$. After multiplying by at most
$2^n$ compositions and summing over $k\le(1-\delta)n$, division by
$(n!)^s$ therefore leaves
\begin{equation}
 \exp\bigl(-s\delta n\log n+O(n)\bigr)
 \label{eq:length-loss}
\end{equation}
for every fixed $\delta\in(0,1)$.

For $k\ge(1-\delta)n$, with $\delta<1/2$, set $D=kx$. Then
\begin{align*}
 \sup_D\e^{-D}(aD^r+Cn)^k
 &=k^{rk}\left[\sup_{x\ge0}
          \e^{-x}\left(ax^r+\frac{Cn}{k^r}\right)\right]^k\\
 &\le k^{rk}\left(a(r/\e)^r+o(1)\right)^k,
\end{align*}
where $o(1)$ is uniform in this range of $k$, since $r>1$.
Stirling's estimate now gives
\begin{equation}
 M_{n,k}\le\e^{o(n)}(ar^r+o(1))^k(k!)^s.
 \label{eq:sharp-type}
\end{equation}
The sum of the numbers of compositions in this range is at most
$\e^{n\mathfrak h(\delta)+o(n)}$, where
$\mathfrak h(\delta)=-\delta\log\delta-(1-\delta)\log(1-\delta)$.
Using $k!\le n!$, the $n$th root of the normalized contribution is at
most
\[
 \e^{\mathfrak h(\delta)}
      \max_{1-\delta\le t\le1}(ar^r)^t+o(1).
\]
First let $n\to\infty$ and then $\delta\downarrow0$.
Together with \eqref{eq:length-loss}, this proves $T_s(\U)\le ar^r$.
All the estimates are upper bounds on positive summands and therefore
remain valid for a restricted subcollection.
\end{proof}

The loss in \eqref{eq:length-loss} is stronger than exponential. The
sharp type comes from compositions whose length is $n-o(n)$; their
entropy has vanishing rate. This is why replacing every composition
count by $2^n$ throughout would destroy the exact constant.

\subsection{A matching one-action lower bound}
Suppose \eqref{eq:strong} holds. We first record the implication
\begin{equation}
 j=o(d_j),\qquad j\log d_j=o(d_j).
 \label{eq:strong-implication}
\end{equation}
Indeed write $d_j=M_jj\log(j+1)$ with $M_j\to\infty$. Since
$\log x/x$ is decreasing for large $x$, or directly from this expression,
\[
 \frac{j\log d_j}{d_j}
 =\frac{\log M_j+\log j+\log\log(j+1)}{M_j\log(j+1)}\longrightarrow0.
\]

Fix $r=s+1$ and put $A=\limsup_j\lambda_j/d_j^r$.
For every $a>A$ when $A<\infty$, the finite exceptions can be included
in a constant $C$ so that \eqref{eq:cost-power-bound} holds.
Lemma~\ref{lem:type-upper} proves $T_s(\U)\le ar^r$; taking $a\downarrow A$
also covers $A=0$.

For the lower bound take any $B>0$ smaller than $A$, or any $B>0$ if
$A=\infty$. There are arbitrarily large $j$ with
$\lambda_j\ge B d_j^r$. Put
$n_j=\lfloor d_j/r\rfloor$ and $m_j=n_j-j$. Equations
\eqref{eq:strong-implication} ensure $m_j\ge0$ eventually,
$m_j/n_j\to1$, and $j\log n_j/n_j\to0$. The one-action term gives
\[
 \frac{u_{n_j}}{(n_j!)^{r-1}}
 \ge\frac{\e^{-d_j}(B d_j^r)^{m_j}}
                 {m_j!(n_j!)^{r-1}}.
\]
Stirling's estimate, $d_j/n_j\to r$, and
$j\log n_j/n_j\to0$ imply
\begin{equation}
 \frac1{n_j}\log\left(
 \frac{\e^{-d_j}(B d_j^r)^{m_j}}
                       {m_j!(n_j!)^{r-1}}\right)
       =\log B+r\log r+o(1).
 \label{eq:lower-type}
\end{equation}
Thus $T_s(\U)\ge Br^r$. Letting $B\uparrow A$ proves
Theorem~\ref{thm:type} for finite positive $A$, and arbitrary $B$ proves
it when $A=\infty$. When $A=0$ the upper bound already gives equality.
No density assumption on the sequence of costs was used: a subsequence
is enough for a limsup.

\subsection{What the type identity does and does not say}
The identity allows isolated large slopes and arbitrarily irregular
costs. In particular, finite changes to amplitudes or slopes do not
change this exact type, although they do change the coefficients and
can change subexponential factors substantially. A type theorem gives
only
$\limsup (u_n/(n!)^s)^{1/n}$; it need not give a limit and never by
itself gives a multiplicative equivalent for $u_n$. Signed reversion
coefficients require a different argument, so no inverse-type identity
is inferred here.

\section{Dense power costs and a dominant-action limit}
\label{sec:condensation}

Assume \eqref{eq:powers}. Let $r=p/\beta$, $s=r-1$, and
$A=a/b^r$. Then $d_j/(j\log(j+1))\to\infty$ and
$\lambda_j/d_j^r\to A$. Theorem~\ref{thm:type} already gives an upper
coefficient-root limit $Ar^r$.

\subsection{Obtaining the lower bound at every degree}
For each large $n$, choose an integer $j_n$ nearest to
$(rn/b)^{1/\beta}$. Since $\beta>1$,
\[
 j_n=o(n),\qquad j_n\log n=o(n),\qquad
 d_{j_n}/n\to r,\qquad \lambda_{j_n}/d_{j_n}^r\to A.
\]
The one-action lower bound with $m=n-j_n$ now applies at every degree
$n$. The calculation in \eqref{eq:lower-type} gives
\[
 \liminf_n\left(\frac{u_n}{(n!)^s}\right)^{1/n}\ge Ar^r.
\]
This proves \eqref{eq:root-limit}. Equivalently,
\begin{equation}
 \log u_n=s\log(n!)+n\log(Ar^r)+o(n).
 \label{eq:log-coefficient}
\end{equation}
It is not the assertion $u_n\sim(Ar^r)^n(n!)^s$.

\subsection{The probability space and the length concentration}
For an ordered composition $\mathbf j=(j_1,\ldots,j_k)$ of $n$, define
\begin{equation}
 \mathbb P_n(\mathbf j)=\frac{1}{u_n k!}
       \e^{-\sum_i d_{j_i}}
       \left(\sum_i\lambda_{j_i}\right)^{k-1}.
 \label{eq:probability}
\end{equation}
Formula~\eqref{eq:ordered-Lagrange} says these probabilities sum to one.
A term with zero weight has zero probability. The variable $K$ is the
number of parts, and $D$ is their total amplitude cost.

For every fixed $\delta>0$, \eqref{eq:length-loss} and the all-degree
lower bound \eqref{eq:log-coefficient} imply
\begin{equation}
 \mathbb P_n\{K\le(1-\delta)n\}
 \le\exp\bigl(-s\delta n\log n+O(n)+o(n)\bigr)\longrightarrow0.
 \label{eq:length-concentration}
\end{equation}
Hence $K/n\to1$. This does not say there is only one nonunit part;
it says the aggregate loss in the number of parts is sublinear.

\subsection{Total cost is concentrated at the unique variational maximum}
The scalar function
\begin{equation}
 x\longmapsto \e^{-x}x^r
 \label{eq:cost-profile}
\end{equation}
has its unique maximum $(r/\e)^r$ at $x=r$. For every fixed
$\varepsilon>0$, its supremum on $|x-r|\ge\varepsilon$ is strictly
smaller. This remains a strict gap after adding a quantity tending
uniformly to zero, as in the proof of \eqref{eq:sharp-type}.

Fix a tolerance $\varepsilon>0$ for $D/n$. Restrict to
$K/n\ge1-\delta$, with $\delta$ sufficiently small. If
$|D/n-r|\ge\varepsilon$, then $|D/K-r|\ge\varepsilon/2$ after possibly
shrinking $\delta$; this follows separately for the upper and lower
inequalities. In the sharp-type proof, the supremum over $x=D/K$ is
therefore restricted away from $r$. It loses a fixed exponential factor.
Use a slope bound with coefficient $a_+>A$ as close to $A$ as needed,
and then take $\delta$ so small that the composition entropy
$\mathfrak h(\delta)$ is smaller than that loss. The contribution has
normalized coefficient-root upper limit strictly below $Ar^r$.
Equation~\eqref{eq:log-coefficient} makes its probability tend to zero.
The excluded range $K/n<1-\delta$ is negligible by
\eqref{eq:length-concentration}. Thus $D/n\to r$ in probability.

For clarity, the two choices are compatible: the continuous function
in \eqref{eq:cost-profile} has a positive gap for the fixed tolerance;
choose $a_+-A$ first small compared with that gap and then choose
$\delta$ small compared with both the gap and the tolerance. The tail
$x\to\infty$ causes no difficulty because \eqref{eq:cost-profile}
tends to zero.

\subsection{A single part carries asymptotically all the cost}
Let $M=\max_i d_{j_i}$. On the event $M\le(1-\varepsilon)D$,
\begin{equation}
 \sum_i d_{j_i}^r\le M^{r-1}D
                     \le(1-\varepsilon)^{r-1}D^r.
 \label{eq:strict-convexity-loss}
\end{equation}
Write $\gamma=(1-\varepsilon)^{r-1}<1$.
Since $\lambda_j\le a_+d_j^r+C$ for every $a_+>A$, the slopes on this
event satisfy $\Lambda\le a_+\gamma D^r+CK$.
The restricted-sum version of Lemma~\ref{lem:type-upper} gives an upper
type $a_+\gamma r^r$. Choose $a_+$ close enough to $A$ that
$a_+\gamma<A$. The all-degree root limit then shows that the probability
of the event tends to zero exponentially. As $D/n\to r>0$, the event
$D=0$ is negligible. Therefore $M/D\to1$ in probability.

It follows that $M/n\to r$. The asymptotic $d_j\sim bj^\beta$ implies
$M/(bJ_{\max}^\beta)\to1$ whenever $J_{\max}\to\infty$: finite initial
indices contribute a bounded cost, and the asymptotic bounds apply
uniformly to all larger indices. Since $M/n\to r$, that divergence
holds in probability. Consequently
\[
 J_{\max}/n^{1/\beta}\to(r/b)^{1/\beta}.
\]
This completes every assertion of Theorem~\ref{thm:condensation}.

\begin{remark}[Meaning of condensation here]
The proved concentration concerns a probability distribution on the
positive Lagrange contributions, not a random solution of the feedback
equation. One action carries almost all amplitude cost and hence almost
all of the power-law slope budget. The theorem gives neither its local
limit law nor the distribution of the remaining small actions. Those
would be needed for a sharp multiplicative coefficient asymptotic.
\end{remark}

\section{A sharper Borel region under strong damping}
\label{sec:sharp-borel}

\begin{proposition}[Damping-controlled inverse continuation]
\label{prop:sharp-borel}
Assume \eqref{eq:strong} and let
$A=\limsup_j\lambda_j/d_j^2<\infty$. If $A>0$, the explicit inverse
Borel series \eqref{eq:borel-series} converges normally on
\begin{equation}
 \mathcal P_A=\left\{\zeta:\ |\zeta|+\Ree\zeta<\frac1{2A}\right\}.
 \label{eq:sharp-parabola}
\end{equation}
If $A=0$, it defines an entire function. In both cases it has at most
exponential growth on a sufficiently thin fixed-width tube around the
negative ray. When $0<A<\infty$, the forward Borel transform has radius
of convergence exactly $1/(4A)$ and has a singularity at the positive
point $1/(4A)$. Also $T_1(\Q)\le4A$; equality and a singularity of the
inverse at the same point are not asserted.
\end{proposition}
\begin{proof}
For every $a>A$ there is a constant $C$ with
$\lambda_j\le a d_j^2+C$. Hence every positive inverse-block frequency
satisfies $b\le aD^2+Ck$ and
$\sqrt b\le\sqrt a\,D+\sqrt{Ck}$. Equation~\eqref{eq:bessel-bound}
now gives
\[
 \e^{-D}|K_{k,b}(\zeta)|
 \le\frac{|\zeta|^k}{k!}
    \e^{-(1-2\sqrt a\,h(\zeta))D}
    \e^{2\sqrt C\,h(\zeta)\sqrt k}.
\]
On compact sets with $2\sqrt a\,h<1$, the first exponential is at most
one. Negative frequencies contribute at most
$\exp(2\sqrt{2\lambda_1}\widetilde h\sqrt k)$ instead. The mass bound
and the factorial in the denominator prove normal convergence. Taking
the union over $a>A$ gives \eqref{eq:sharp-parabola}; for $A=0$, every
compact subset of the plane is permitted by taking $a$ small enough.

On a fixed tube, $h$ is bounded and $\widetilde h=O(\sqrt{1+|\zeta|})$.
Young's inequality absorbs the remaining square-root terms into
$C_1k+C_2(1+|\zeta|)$, and the sum over $k$ is exponentially bounded.
This argument is simpler than \eqref{eq:tube-log}, because the strong
cost estimate has removed $k\log k$ from the positive-frequency bound.

Theorem~\ref{thm:type} at $s=1$ gives $T_1(\U)=4A$, hence the stated
forward Borel radius. A positive-coefficient power series of finite
radius $R$ cannot be analytic at $R$: if it were, take a real $x<R$
close enough to $R$ that its Taylor expansion about $x$ converges at
some $y>R$. Its Taylor coefficients at $x$ are nonnegative. Interchanging
the resulting nonnegative double sum would show that the original
power series converges at $y$, a contradiction. This classical
positive-coefficient argument proves the forward singularity.
Finally, \eqref{eq:sharp-parabola} contains every disc of radius less
than $1/(4A)$, so the inverse Borel radius is at least $1/(4A)$.
\end{proof}

The positive intercept of the guaranteed inverse parabola is therefore
exactly the forward Borel radius. This is a useful alignment of two
independent calculations, but it does not prove that the boundary of
\eqref{eq:sharp-parabola} is a natural boundary, or even that every
point on it is singular for the inverse. Normal convergence regions
and maximal continuation domains must not be identified.

\section{Phase diagrams and explicit examples}
\label{sec:examples}

\subsection{Power-law slopes and stretched-exponential amplitudes}
Suppose $d_j\sim b j^\beta$ and $\lambda_j\sim a j^p$, with
$a,b,\beta,p>0$, and keep the normalization at $j=1$. Put
$B=\max(1,\beta)$. The sharp classification reads
\begin{center}
\begin{tabular}{@{}p{0.54\linewidth}p{0.38\linewidth}@{}}
\toprule
Property & Necessary and sufficient condition\\
\midrule
Ordinary convergence & $p\le B$\\[3pt]
Gevrey--$s$, for fixed $s>0$ & $p\le(s+1)B$\\[3pt]
Fine summability in direction $\pi$ & $p\le2B$\\[3pt]
Least Gevrey order when $p>B$ & $s_*=p/B-1$, attained\\
\bottomrule
\end{tabular}
\end{center}
At $\beta\le1$, the logarithm in $j\log(j+1)$ makes every equality
case in the Gevrey row permissible. For $\beta>1$ the damping cost
dominates that logarithm. In the divergent range $p>\beta>1$,
\begin{equation}
 T_{p/\beta-1}(\U)
 =a\left(\frac{p}{b\beta}\right)^{p/\beta},
 \qquad
 J_{\max}\sim_{\mathbb P}
       \left(\frac{p}{b\beta}n\right)^{1/\beta}.
 \label{eq:power-summary}
\end{equation}
The notation $\sim_{\mathbb P}$ denotes a ratio tending to one in
probability under \eqref{eq:probability}.

\begin{example}[Quadratic costs and cubic slopes]
Let $c_j=2^{-j^2}$ and $\lambda_j=j^3$ for $j\ge2$, with $c_1=1$ and
$\lambda_1=1$. Here $\beta=2$, $b=\log2$, $p=3$. The formal solution
is divergent, has least Gevrey order $1/2$, and satisfies
\[
 \lim_n\left(\frac{u_n}{\sqrt{n!}}\right)^{1/n}
       =\left(\frac{3}{2\log2}\right)^{3/2}.
\]
Both series are finely summable in the negative direction, since they
are in particular Gevrey--1. Their ordinary Borel transforms are entire
in this case: for the forward series this follows from its order-$1/2$
bound, and for the inverse from Proposition~\ref{prop:sharp-borel} with
$A=0$. Entire does not imply exponential growth in every direction.
\end{example}

\begin{example}[A nonzero critical Borel scale]
For $j\ge2$, take $c_j=\exp(-b j^\beta)$ and
$\lambda_j=a j^{2\beta}$, where $\beta>1$ and $a,b>0$.
The least Gevrey order is one, with exact forward type $4a/b^2$.
The forward Borel transform has its first positive singularity at
$\zeta=b^2/(4a)$. The inverse has the guaranteed continuation region
\[
 |\zeta|+\Ree\zeta<\frac{b^2}{2a}.
\]
It is finely summable on the negative ray. Nothing in this argument
classifies the singularity at the positive point or its Stokes data.
\end{example}

\subsection{Logarithmic boundary factors}
The budgets also settle endpoints not visible in a pure power diagram.
Suppose, for $j\to\infty$,
\[
 d_j\sim b j^\beta(\log j)^\gamma,\qquad
 \lambda_j\sim a j^p(\log j)^\eta,
 \quad a,b,\beta,p>0,\quad\gamma,\eta\in\R.
\]
Order pairs lexicographically. Let
\[
 (B_0,G_0)=\max_{\rm lex}\{(1,0),(\beta,\gamma)\},\qquad
 (B_1,G_1)=\max_{\rm lex}\{(1,1),(\beta,\gamma)\}.
\]
Comparison of positive monomials gives the complete tests
\begin{align}
 \U\text{ converges}
 &\iff (p,\eta)\le_{\rm lex}(B_0,G_0),\label{eq:log-analytic}\\
 \U,\Q\in\Gev_s
 &\iff \left(\frac p{s+1},\frac\eta{s+1}\right)
                       \le_{\rm lex}(B_1,G_1).
 \label{eq:log-gevrey}
\end{align}
The second test at $s=1$ is also the fine-summability test. For example,
when $d_j\sim b j^\beta$ with $\beta>1$ and $p>\beta$, the infimal
positive order is $p/\beta-1$. At that order a factor
$(\log j)^\eta$ in the slopes is permitted exactly when $\eta\le0$.
Thus an infimal Gevrey order need not be attained in the general
logarithmic setting. The main criterion records that distinction.

\subsection{Superpolynomial slopes can still be summable}
\begin{example}[Exponential slopes, exact Gevrey--1]
For $j\ge2$ set
\[
 c_j=\exp(-\e^{j/2}),\qquad \lambda_j=\e^j,
\]
and set $c_1=1$, $\lambda_1=1$. Then $d_j=\e^{j/2}$ and
$\lambda_j=d_j^2$. The convergence criterion fails, but the least Gevrey
order is exactly one and $T_1(\U)=4$. The forward Borel radius is
$1/4$; the inverse has the guaranteed region
$|\zeta|+\Ree\zeta<1/2$; both series are finely summable in direction
$\pi$. The geometric spacing of these costs does not meet the power-cost
hypothesis of Theorem~\ref{thm:condensation}, so no all-degree root limit
is claimed by that theorem.
\end{example}

\begin{example}[Damping sufficient for ordinary convergence]
Let $\lambda_j\ge0$ be any sequence, however fast it grows, and take
$c_j=\exp(-\lambda_j)$ for $j\ge2$, with $c_1=1$.
Then $d_j=\lambda_j$ for $j\ge2$, so
$\lambda_j/(j+d_j)\le1$. The formal solution is genuinely convergent.
The amplitude cost can therefore restore ordinary analyticity, not just
a weaker factorial regularity. Finitely many initial slopes do not
affect this conclusion.
\end{example}

\begin{example}[A divergent series in every positive Gevrey class]
Let $c_j=1$ and $\lambda_j=j\log(j+1)$. Then
$\lambda_j/j\to\infty$, so the solution diverges. For every $s>0$,
$\lambda_j^{1/(s+1)}/(j\log(j+1))$ is bounded. Thus the solution and
its inverse belong to every positive Gevrey class and are finely
summable in direction $\pi$. This is an illustration of the distinction
between the two budgets, not a new unit-amplitude theorem.
\end{example}

\section{Why the hypotheses and the direction matter}
\label{sec:boundaries}

\subsection{Exponential amplitude control is a substantive assumption}
The joint fine-summability equivalence is not a theorem for arbitrary
positive amplitudes. Take $\lambda_j=0$ for all $j$ and
\[
 c_j=j!\bigl(2+(-1)^j\bigr),\qquad j\ge1.
\]
Then $c_1=1$ and the equation is simply $\U(q)=\sum c_jq^j$.
The series is Gevrey--1, but its integrated Borel transform is
\begin{equation}
 \B\U(\zeta)=\frac{2\zeta}{1-\zeta}
                         -\frac{\zeta}{1+\zeta}.
 \label{eq:amplitude-counterexample}
\end{equation}
The pole at $\zeta=-1$ prevents fine summability in the negative
direction. The amplitudes are not exponentially bounded, so this does
not contradict Theorem~\ref{thm:fine}. It shows exactly where the
geometric block mass used in the proof can fail.

\subsection{Positivity is not just a convenient estimate}
The one-action lower bound, the exact-type lower bound, and the
probability interpretation use nonnegative Lagrange terms. Signed or
complex amplitudes can cancel those terms. A theorem for such models
would need a noncancellation condition or different invariants.
Similarly, negative or complex slopes destroy the common damping
half-plane used to sum the literal kernel and to control Bessel
frequencies. The present necessity-and-sufficiency statements should
not be extended to those settings by taking absolute values in their
conclusions.

\subsection{Why positive-ray summation fails for divergent positive series}
The asymmetry between positive and negative directions is intrinsic.
Suppose $\U$ has nonnegative coefficients, is divergent, and is
Gevrey--1. If its Borel transform has finite radius, the
positive-coefficient argument in Proposition~\ref{prop:sharp-borel}
puts a singularity on the positive ray. If its Borel transform is entire
and had an exponential bound $\B\U(t)\le C\e^{Ht}$ for $t\ge0$,
positivity would give
\[
 \frac{u_n}{n!}\le\frac{C\e^{Ht}}{t^n}.
\]
Taking $t=n/H$ (and increasing $H$ to be positive if necessary) and
using Stirling's estimate would yield
$u_n\le C'\sqrt n\,H^n$, forcing ordinary convergence. This is a
contradiction. Thus a divergent positive solution cannot be finely
Borel summable in the positive direction, even in cases where its
Borel transform is entire.

\subsection{Remaining distinctions}
The literal inverse is analytic in a small open left half-disc and has
continuous boundary values supplied by its blocks. The uniform
Gevrey--$s$ remainder theorem is proved for that inverse. For $s=1$,
fine summation gives a canonical value in its summation class and an
inverse forward sum. These statements do not imply a common open arc
of summation directions, resurgent continuation along arbitrary paths,
a sharp least-truncation error, or a Stokes formula.

For $s\ne1$ the factorial remainder theorem does not, by itself, prove
summability with a moment kernel of order $s$. In particular, the
all-slope continuous ray transform in \eqref{eq:ray-transform} is not
to be relabeled as a holomorphic Borel transform when its germ fails
to exist. The separation of these assertions is part of the theorem.

\section{Exact checks, numerical diagnostics, and reproducibility}
\label{sec:verification}

\subsection{Independent finite identities}
The package's primary verification script, \texttt{code/verify.py},
uses exact rational arithmetic from Python's standard library. Its
recorded run completed \textbf{1,353 assertions} at coefficient order
12 across twelve models. The models combine several integer-power
slope sequences with unit, geometric, quadratic-cost, and irregular
positive rational amplitudes. The checks compare two independent
forward constructions: the exponential coefficient recurrence and the
ordered Lagrange formula. The inverse is constructed from the signed
exponential blocks, after which both compositions $\U\circ\Q$ and
$\Q\circ\U$ are checked coefficientwise.

For reference, the recurrence used independently of Lagrange inversion
is as follows. Write
\[
 E_j(q)=\exp(\lambda_j\U(q))=\sum_{k\ge0}e_{j,k}q^k.
\]
Differentiation and coefficient comparison give
\begin{equation}
 e_{j,0}=1,\qquad
 e_{j,k}=\frac{\lambda_j}{k}\sum_{m=1}^k m u_m e_{j,k-m},
 \qquad
 u_n=\sum_{j=1}^n c_j e_{j,n-j}.
 \label{eq:algorithm}
\end{equation}
Every $e_{j,n-j}$ on the right uses only previously computed $u_m$.
No nonlinear numerical root finder is needed. The other exact checks
cover the one-action lower bound and the block mass bound through block
100. The certificate below uses explicit rational exponential
enclosures together with checks on its domain and final width. Exact finite checks test identities and implementation consistency;
they are not proofs of the asymptotic statements.

\subsection{A rational certificate for a literal inverse value}
For the cubic-slope, quadratic-cost model
\[
 c_1=1,\quad c_j=2^{-j^2}\ (j\ge2),\qquad \lambda_j=j^3,
\]
the script encloses the first ten inverse blocks at $u=-1/100$ using
rational exponential Taylor bounds and adds the geometric tail
\eqref{eq:block-certificate}. It produces the following outward-rounded
decimal enclosure:
\begin{equation}
 -0.010106452308614
 < Q(-0.01) <
 -0.010106452306282.
 \label{eq:certificate-value}
\end{equation}
The file \path{data/exact_checks.json} records the exact rational
endpoints. Their interval has width less than $2.331\times10^{-12}$;
the displayed decimals are rounded outward. This certifies the literal
implicit branch, rather than agreement between two floating-point
implementations. The tail bound ignores cancellation and is not optimized.

An independent interval evaluation of the literal kernel, using its first
15 actions and a geometric tail, confirms opposite residual signs at
these rational endpoints. On $|q|\le1/50$, its derivative is bounded below by
\[
 1-\frac1{100}-\left((1-1/50)^{-2}-1\right)>0.
\]
Thus the same interval contains a unique literal root by a second route.
The endpoint sign enclosures and derivative bound are recorded alongside
the block certificate.


\subsection{Coefficient-root diagnostics}
The separate script \texttt{code/diagnostics.py} evaluates
\eqref{eq:algorithm} in logarithmic floating-point arithmetic, using
NumPy and SciPy. It was first compared with exact rational coefficients
through degree 20; the largest discrepancy between logarithms was
$4.27\times10^{-14}$ after upward rounding. The subsequent calculations
through degree 512 are diagnostics, not interval certificates.

The table uses $c_j=2^{-j^2}$ for $j\ge2$, $c_1=1$,
$\lambda_j=j^p$, and
\[
 R_n=\frac{(u_n/(n!)^{p/2-1})^{1/n}}
                {(p/(2\log2))^{p/2}}.
\]
Theorem~\ref{thm:condensation} gives $R_n\to1$.
\begin{center}
\begin{tabular}{@{}rccc@{}}
\toprule
$n$ & $p=3$ & $p=4$ & $p=5$\\
\midrule
64  & 0.592248 & 0.305540 & 0.149857\\
128 & 0.658251 & 0.384842 & 0.213325\\
256 & 0.719314 & 0.468507 & 0.290177\\
512 & 0.773542 & 0.551826 & 0.376472\\
\bottomrule
\end{tabular}
\end{center}
Convergence is visibly slow at these degrees. The discrepancy is not an
error bar for the theorem and is not evidence for a multiplicative
coefficient equivalent. The recorded diagnostic file additionally
contains large-degree one-action lower bounds, which are inexpensive to
compute but are not the full coefficients. Keeping those two outputs
separate prevents a lower-bound calculation from being mistaken for an
exact coefficient evaluation.

\subsection{Reproduction commands and dependencies}
The article is self-contained LaTeX and uses standard packages. From
the package directory, run
\begin{quote}\small\ttfamily\raggedright
python code/verify.py --order 12\\
python code/diagnostics.py\\
pdflatex -interaction=nonstopmode -halt-on-error article.tex\\
pdflatex -interaction=nonstopmode -halt-on-error article.tex\\
pdflatex -interaction=nonstopmode -halt-on-error article.tex
\end{quote}
The supplied \texttt{build.sh} runs the three document passes. The exact
script needs no third-party Python dependency; only the diagnostic
script needs NumPy and SciPy. The scripts make no network requests.
The package contains the scripts, their recorded outputs, source notes,
and build instructions. No source or branch in ProveIt was modified.

\section{Formalization targets and proof dependencies}
\label{sec:formalization}

The repository already has formal infrastructure for asymptotic scales,
flatness, well-based supports, and differential blocks \cite{project}.
Those declarations do not automatically formalize this article's
infinite analytic constructions. A useful implementation should separate
four layers and record the exact theorem supplied by each.

\subsection{Finite coefficient algebra}
The first layer should define weighted compositions and the finite
Lagrange coefficients over a commutative algebra, prove the recurrence
\eqref{eq:algorithm}, and identify the signed inverse blocks. Rational
specializations can then produce kernel-checked finite identities.
Positivity and the one-action injection give a particularly small
initial target. These lemmas have no dependence on Borel summation,
asymptotic estimates, or numerical special functions.

\subsection{Real estimates and the factorial budget}
The next layer should establish the block mass estimate, the cost-moment
bound, factorial inequalities, and the uniform exponential Taylor
remainder. The compact statement of Lemma~\ref{lem:factorial-budget}
would be a reusable library result. A theorem about a family of finite
moments must be distinguished from a theorem about a function remainder;
the geometric tail and uniform-limit arguments connect them.

\subsection{Analytic realization and summation}
The analytic layer needs uniform convergence on closed half-discs,
normal convergence in open domains, the fixed-parameter inverse theorem,
and the beta-measure kernel identity. The two Laplace exchanges have
different hypotheses: absolute Taylor integration for one frequency,
and an oscillatory absolute-integral bound for infinitely many blocks.
Formalizing them as separate lemmas would guard against a common invalid
exchange. A general fine-summability algebra should then package the
convolution and inverse theorem in Appendix~\ref{app:inverse}.

\subsection{Sharp asymptotics and probability}
The final layer contains the entropy estimate, the two length ranges in
Lemma~\ref{lem:type-upper}, and the conversion of exponential losses
into probability concentration. The statement must distinguish a limsup
type, an all-degree root limit, and a multiplicative equivalent.
The current argument establishes the first two in their respective
hypotheses, not the third. It should likewise distinguish concentration
of the largest cost from any unproved local central limit theorem.
No Lean files are supplied or claimed to have been checked in this
package; this section is an implementation roadmap.

\section{Further research questions}
\label{sec:research}

The following are proposed next tasks for this weighted model. They are
not blanket assertions that every analogous question is open throughout
the literature; comparison with general type-preservation and summation
theorems should precede a publication-priority claim.

\subsection{Sharp signed inverse coefficients}
Determine the exact type and, more ambitiously, signed coefficient
asymptotics for $\Q$ under strong damping. The parabolic construction
already gives $T_1(\Q)\le T_1(\U)$ in the finite positive-type case.
A general exact-type-preservation theorem, if applicable with this
normalization, should be used rather than rediscovered. Even such a
theorem would not settle the signs, cancellation scale, or multiplicative
prefactor of the inverse coefficients.

\subsection{A full multiplicative forward asymptotic}
For $d_j=bj^\beta$ and $\lambda_j=aj^p$, refine
\eqref{eq:log-coefficient} to a multiplicative equivalent. The dominant
large action has size $n^{1/\beta}$, but small nonunit parts, the cost of
using fewer than $n$ parts, and lattice rounding all contribute below
order $n$. Determine how many small actions must be retained to obtain
relative error $o(1)$, and whether a finite-core threshold depends on
both $p$ and $\beta$ rather than on their ratio alone.

\subsection{Local fluctuations and a conditional ensemble}
Prove a local limit theorem for $J_{\max}$ and a fluctuation theorem for
$K$ under \eqref{eq:probability}. A formal quadratic expansion of the
cost profile suggests a first scale, but it does not account for the
constraint that total action is exactly $n$. The first objective is a
correct centering, including the drift from the sea of small actions;
a Gaussian formula with the wrong center is not a local limit theorem.

\subsection{Lacunary cost spectra and log-periodic behavior}
For geometrically spaced costs, such as $d_j=\e^{\theta j}$, the type
proof obtains a sharp limsup by choosing degrees matched to individual
costs. Determine the lower root envelope at unmatched degrees, whether
several expensive actions can compensate for a gap, and whether the
normalized logarithms have periodic limit profiles in $\log n$.
The dense power-cost theorem does not decide this question.

\subsection{Actual inverse Borel singularities}
In the finite positive-type regime, determine the maximal continuation
domain of $B_Q$, beginning with the positive intercept of
\eqref{eq:sharp-parabola}. Does the first forward singularity survive
reversion at the same Borel location? Determine its local form and any
lateral summation discontinuity in specific power-cost models. Normal
convergence of a chosen kernel representation gives only one side of
this problem.

\subsection{Angular and generalized moment summation}
Strengthen fine negative-ray summability to an open arc of directions
where possible, and identify obstructions where it fails. For $s\ne1$,
construct moment kernels appropriate to $(n!)^s$ or
$\Gamma(1+sn)$, keeping the exact conversion of types visible.
The target should identify the sum with the literal inverse, not merely
prove the existence of some function with the same asymptotic expansion.
Uniform Gevrey remainders are a starting point, not the entire summation
argument.

\subsection{Amplitudes beyond exponential growth}
The counterexample \eqref{eq:amplitude-counterexample} shows that slope
control alone cannot handle rapidly growing positive amplitudes. Seek a
joint criterion in terms of the amplitude generating function's Borel
continuation and the cost of its interaction with the slopes. Even the
zero-slope model forces directional singularity information into such
a criterion. This suggests separating an amplitude summation class from
a feedback stability condition.

\subsection{Complex phases and vector feedback}
For signed amplitudes or complex slopes, identify usable
noncancellation or sectorial-damping conditions. For a vector equation,
replace the scalar exponential feedback by linear forms in several
unknowns and ask whether a matrix or convex-geometric budget controls
regularity. A countable-tail condition must be explicit; a finite-dimensional
implicit-function theorem cannot control an unbounded family of
coefficients by itself.

\subsection{General actions and missing action one}
Replace $q^j$ by $q^{a_j}$ for integer or well-based real actions.
Determine which version of the budget survives after replacing the
index $j$ by action size. The present lower bound uses arbitrarily many
action-one factors and the present inverse uses an ordinary nonzero
linear coefficient. When the smallest action is different, ramification
or a different formal inverse object must be specified before an
analytic statement can be formulated.

\subsection{Certified algorithms with adaptive costs}
The universal factor $8^k$ is simple but wasteful. Develop certified
weighted mass generating functions and choose adaptively between exact
block truncation, Taylor truncation, and Borel quadrature. The amplitude
cost is available before a frequency is evaluated and can guide
pruning of negligible compositions. A useful deliverable would combine
a formally checked finite certificate with an explicit analytic tail
bound, rather than rely on floating cancellation at high order.

\section{Conclusion}
The positive weighted feedback model admits a sharp and explicit answer
to the amplitude--slope question. Ordinary analyticity is governed by
$j+d_j$, while every positive Gevrey order is governed by
$j\log(j+1)+d_j$. The same latter budget controls an actual uniform
left-side inverse remainder; at order one it also characterizes fine
negative-direction summability and identifies the summed equation.

Strong damping produces a sharper theorem: an exact type controlled by
$\lambda_j/d_j^{s+1}$, and, for dense power costs, an all-degree root
limit and concentration of one dominant action. These are model-specific
results with complete arguments and reproducible finite checks. They do
not rely on accepting an unreviewed repository draft as a proof, and
they leave the multiplicative asymptotics, signed inverse behavior,
maximal Borel continuation, and stronger summation questions explicitly
separate.

\appendix
\section{Compositional inversion in the fine-summable class}
\label{app:inverse}

This appendix proves the nonlinear result used in
Theorem~\ref{thm:fine}. It is included because a theorem stated for an
open angular sector cannot silently be applied to a fixed-width tube.
The proof specializes the classical convolution method described in
\cite{sauzin}.

\subsection{The Borel minor at infinity}
For $\varphi(z)=\sum_{n\ge1}\varphi_n z^{-n}$ define
\[
 \mathscr B\varphi(\xi)=\sum_{n\ge1}
                \frac{\varphi_n}{(n-1)!}\xi^{n-1}.
\]
Let $\Omega^+_\delta=-\Omega^-_\delta$, which is convex. Suppose a minor
$\widehat\varphi$ is holomorphic there and
$|\widehat\varphi(\xi)|\le M\e^{H|\xi|}$. Products of formal series
correspond to convolution along the segment from zero to $\xi$.
The convolution simplex has volume $1/(k-1)!$, and convexity keeps
all its arguments in the tube. Thus
\begin{equation}
 |\widehat\varphi^{*k}(\xi)|
 \le M^k\frac{|\xi|^{k-1}}{(k-1)!}\e^{H|\xi|}.
 \label{eq:convolution-bound}
\end{equation}
The sum of the lengths of the collinear arguments is $|\xi|$, which
explains the single exponential on the right.

If $F(w)=\sum F_kw^k$ is convergent near zero, with
$|F_k|\le C R^{-k}$, \eqref{eq:convolution-bound} proves normal convergence
of the minor of $F(\varphi)-F(0)$ and an exponential bound. In particular,
reciprocals of formal series with nonzero constant term are permitted.
Multiplication by $z$, after removal of the new constant term, is
differentiation of the minor and preserves its exponential bound on a
narrower tube by Cauchy's estimate. Division by $z$ is integration.
Translation $z\mapsto z+c$ multiplies the minor by $\e^{-c\xi}$.
All these operations preserve this fine-summable class.

Laplace summation
\[
 \mathcal S\varphi(z)=\int_0^\infty
                  \e^{-z\xi}\widehat\varphi(\xi)\,\dd\xi,
 \qquad \Ree z>H,
\]
commutes with the stated operations. For products this is the
absolutely convergent convolution integral and Fubini's theorem; for
analytic scalar substitution it follows from the same exponential
majorant. Multiplication by $z$ follows by integration by parts, and
translation follows directly from its multiplier. These arguments also
give the estimates $|\mathcal S\varphi(z)|\le M/(\Ree z-H)$ and
$|\partial_z\mathcal S\varphi(z)|\le M/(\Ree z-H)^2$.

\subsection{An explicit inverse minor}
\begin{lemma}[Fine inversion at infinity]
\label{lem:fine-inverse-infinity}
Let $f(z)=z+c+\varphi(z)$, where $\varphi=O(z^{-1})$ has a minor
holomorphic with an exponential bound on some $\Omega^+_\delta$.
Its formal inverse has the form $g(z)=z-c+\chi(z)$ with $\chi=O(z^{-1})$
in the same class, possibly on a narrower tube. Their Laplace sums are
analytic inverses on sufficiently far right half-planes and their
corresponding images.
\end{lemma}
\begin{proof}
First take $c=0$. Formal Lagrange inversion gives
\begin{equation}
 \chi(z)=\sum_{k\ge1}\frac{(-1)^k}{k!}
                       \partial_z^{k-1}\bigl(\varphi(z)^k\bigr).
 \label{eq:inverse-operator}
\end{equation}
For example, apply coefficient Lagrange inversion to
$w=-t\varphi(z+w)$ and then put $t=1$; the $z^{-1}$-order of the
$k$th summand tends to infinity. Differentiation corresponds to
multiplication of a Borel minor by $-\xi$, so the proposed inverse minor
is
\begin{equation}
 \widehat\chi(\xi)=-\sum_{k\ge1}\frac{\xi^{k-1}}{k!}
                                     \widehat\varphi^{*k}(\xi).
 \label{eq:inverse-minor}
\end{equation}
Equation~\eqref{eq:convolution-bound} bounds its absolute value by
\[
 \e^{H|\xi|}\sum_{k\ge1}
       \frac{M^k|\xi|^{2k-2}}{k!(k-1)!}
 \le M\e^{(H+2\sqrt M)|\xi|}.
\]
The last inequality follows, for $t\ge0$, from
$\sum_{\ell\ge0}t^{2\ell}/(\ell!(\ell+1)!)\le\e^{2t}$,
using $\binom{2\ell}{\ell}\le4^\ell$.
This also proves normal convergence and identifies the formal
coefficients.

It remains to identify the analytic sum, rather than just its formal
inverse coefficients. On a far right half-plane, the Laplace sum
$\phi(z)=\mathcal S\varphi(z)$ is uniformly small and has uniformly
small derivative. The equation $w=-t\phi(z+w)$ has a holomorphic
solution for $|t|$ in a disc of radius greater than one: choose a fixed
small disc for $w$ and move the half-plane sufficiently far right to
make the map a strict contraction there. Classical analytic Lagrange
inversion yields \eqref{eq:inverse-operator} with $\phi$ in place of
$\varphi$, evaluated at $t=1$.

The Laplace integral of the absolute $k$th minor is at most
\[
 \frac{M^k(2k-2)!}{k!(k-1)!(\Ree z-H)^{2k-1}},
\]
which is summable when $\Ree z-H>2\sqrt M$.
Thus summation of \eqref{eq:inverse-minor} commutes with its Laplace
integral and gives exactly that analytic Lagrange series.
The sum $z+\mathcal S\chi(z)$ is the analytic inverse. Injectivity on
a far right half-plane also follows by integrating the derivative
bound $|\phi'|<1/2$ along line segments in that convex domain.
For general $c$, invert $f_0(z)=z+\varphi(z)$ first and then translate:
$f^{-1}(z)=f_0^{-1}(z-c)$. Translation preserves all the established
properties.
\end{proof}

\subsection{Returning to series tangent to the identity at zero}
Let $Q(u)=u+O(u^2)$ be a fine-summable formal series in direction $\pi$.
Conjugate it by $T(z)=-1/z$:
\[
 F(z)=T\circ Q\circ T(z)=-\frac1{Q(-1/z)}=z+c+\varphi(z).
\]
Changing the sign of the Borel variable moves the negative tube to the
positive one. The minor at infinity differs from the integrated Borel
transform only by differentiation, and the reciprocal and multiplication
operations used to form $F$ are permitted by the preceding algebra.
Lemma~\ref{lem:fine-inverse-infinity} constructs $G=F^{-1}$ in the same
class and identifies its sum. Conjugating back gives
\[
 T\circ G\circ T=Q^{\inv},
\]
which is finely summable in direction $\pi$. Summation commutes with the
conjugation operations, so the sums are inverse functions on their
small common domains. This is the exact closure result used in the
main text; no angular-sector hypothesis has been inserted.

\section{Notation and the logical dependency map}
\label{app:notation}

\begin{center}
\begin{tabular}{@{}p{0.22\linewidth}p{0.69\linewidth}@{}}
\toprule
Symbol & Meaning\\
\midrule
$c_j$, $d_j$ & Positive amplitude and its cost $-\log c_j$\\
$\lambda_j$ & Nonnegative feedback slope at action $j$\\
$u_n$, $v_n$ & Forward and inverse ordinary coefficients\\
$L_j$, $H_j$ & $j\log(j+1)$ and $j\log(j+1)+d_j$\\
$s$, $r$ & Gevrey order $s>0$ and $r=s+1$\\
$T_s$ & Type using the normalization $(n!)^s$\\
$k$, $D$, $b$ & Inverse block number, composition cost, frequency\\
$K$, $D$, $J_{\max}$ & Lagrange-ensemble length, cost, largest action\\
$\B$, $\mathscr B$ & Integrated Borel transform at zero; minor at infinity\\
$h$, $\widetilde h$ & Absolute real and imaginary parts of $\sqrt\zeta$\\
\bottomrule
\end{tabular}
\end{center}
The repeated letter $D$ always denotes a sum of amplitude costs; its
indexing composition is clear from the context. The inverse block index
$k$ is not the length $\ell$ of its internal composition. The Lagrange
length $K$ belongs instead to the forward probability ensemble.

The proof dependencies can be summarized without identifying unlike
statements. Finite Lagrange algebra gives positivity and the necessary
budget. Weighted inverse blocks, the moment estimate, and the factorial
lemma give sufficient Gevrey bounds and uniform inverse remainders.
Bessel kernels plus the weighted tube estimate give fine inverse
summation, and the nonlinear lemma transfers it to the forward series.
Separately, a positive variational upper bound and a one-action lower
bound give the exact strong-damping type. Dense power costs make the
lower bound available at every degree and then turn strict variational
losses into probability concentration. None of these arrows uses the
finite computational checks as a substitute for an infinite proof.

\clearpage
\begin{thebibliography}{9}
\bibitem{project}
Vladimir Reshetnikov, \emph{ProveIt: Transseries project and document
inventory}, repository snapshot
\texttt{ebd8344bca77d8352cd745b7df374618a290a029}, inspected 29 September 2026.
\url{https://github.com/VladimirReshetnikov/ProveIt/tree/ebd8344bca77d8352cd745b7df374618a290a029/Analysis/Transseries}.
The project README and
\texttt{docs/series-and-transseries/README.md} supply the provenance and
formalization-status context.

\bibitem{regularity}
\emph{A Sharp Regularity Classification for Countable Exponential-Feedback
Transseries}, research draft in ProveIt, same snapshot, 2026.
\url{https://github.com/VladimirReshetnikov/ProveIt/blob/ebd8344bca77d8352cd745b7df374618a290a029/Analysis/Transseries/docs/series-and-transseries/Exponential_Feedback_Regularity_Classification/article.tex}.
The amplitude question is at source lines 1470--1483; other nearby
questions concern exact inverse type and uniform remainders.

\bibitem{negative}
\emph{Negative-Ray Summation of Countable Exponential-Feedback Transseries},
research draft in ProveIt, same snapshot, 2026.
\url{https://github.com/VladimirReshetnikov/ProveIt/blob/ebd8344bca77d8352cd745b7df374618a290a029/Analysis/Transseries/docs/series-and-transseries/Negative_Ray_Summation_Exponential_Feedback/article.tex}.
The weighted-amplitude question and the distinction between fine and
angular summation are the immediate analytic starting points.

\bibitem{gessel}
Ira M. Gessel, \emph{Lagrange inversion}, Journal of Combinatorial Theory,
Series A \textbf{144} (2016), 212--249.
\url{https://doi.org/10.1016/j.jcta.2016.06.018}.
Preprint: \url{https://arxiv.org/abs/1609.05988}.

\bibitem{sauzin}
David Sauzin, \emph{Introduction to 1-summability and resurgence},
arXiv:1405.0356, 2014.
\url{https://arxiv.org/abs/1405.0356}.
Sections 7 and 9 distinguish fine summation from summation in an arc of
directions; the convolution framework provides the background for the
specialized nonlinear proof in Appendix~\ref{app:inverse}.

\bibitem{dlmf}
NIST Digital Library of Mathematical Functions,
\emph{Bessel functions}, Sections 10.9 and 10.25, accessed 29 September 2026.
\url{https://dlmf.nist.gov/10.9} and
\url{https://dlmf.nist.gov/10.25}.
The Poisson integral and modified Bessel power series underlie the
normalization proved directly in \eqref{eq:bessel-integral}.
\end{thebibliography}
\end{document}
